%%%%%%%%%%%%%%%%%%%%%%%%%%%%%%%%%%%%%%%%%%%%%%%%%%
% tcbset 
% history
% v1.0 : 27 Feb, 2026
% v2.0 : 17 Apr, 2026
% update width for beamer
%%%%%%%%%%%%%%%%%%%%%%%%%%%%%%%%%%%%%%%%%%%%%%%%%%


% Blue: 信頼感のあるスレートブルー系
\definecolor{mod-blue-bg}{HTML}{F0F4F8}
\definecolor{mod-blue-ln}{HTML}{476683}
\definecolor{ProbBlue}{HTML}{E2EAF4}
\definecolor{ProbBorder}{HTML}{3867D6}

% Green: 自然で落ち着いたセージグリーン系
\definecolor{mod-green-bg}{HTML}{F1F4F0}
\definecolor{mod-green-ln}{HTML}{557555}
\definecolor{ExamGreen}{HTML}{E8F2E9}
\definecolor{ExamBorder}{HTML}{2D6A4F}

% Gray/Neutral: 洗練されたウォームグレー・サンド系
\definecolor{mod-gray-bg}{HTML}{F7F7F7}
\definecolor{mod-gray-ln}{HTML}{5F6368}

% Accent/Alert: 重要な定義など（sbox0等に推奨）
\definecolor{mod-accent-bg}{HTML}{FFF9E6} % 非常に薄い琥珀色
\definecolor{mod-accent-ln}{HTML}{D97706} % 落ち着いたオレンジ

%\definecolor{mod-blue-bg}{HTML}{F4F7F9}
%\definecolor{mod-blue-ln}{HTML}{4A6984}
%\definecolor{mod-green-bg}{HTML}{F2F5F1}
%\definecolor{mod-green-ln}{HTML}{5E7D5E}
%\definecolor{mod-gray-bg}{HTML}{F8F8F7}
%\definecolor{mod-gray-ln}{HTML}{50504F}
%\definecolor{ProbBlue}{HTML}{E8F0FE}
%\definecolor{ProbBorder}{HTML}{4285F4}
%\definecolor{ExamGreen}{HTML}{EAF5EA}
%\definecolor{ExamBorder}{HTML}{2E7D32}

\tcbset{
    modern-base/.style={
        %breakable,
        enhanced,
        frame hidden,
        %sharp corners,
        boxsep=5pt,
        left=3pt,
        right=3pt,
        leftrule=0pt,
        toprule=0pt,
        rightrule=0pt,
        bottomrule=0pt,
        before upper={\setlength{\parindent}{1\zw}},
        width=0.9\linewidth,
        center
    }
}

%% --- 各ボックスの再定義 ---

% dbox0: 点線・白背景（極めてシンプル）
\NewTColorBox{dbox0}{}{
    modern-base,
    borderline={0.3pt}{0pt}{black, dashed},
    colback=white
}

% dbox1: 点線・薄い青背景（注釈や導出過程向け）
\NewTColorBox{dbox1}{}{
    modern-base,
    borderline={0.3pt}{0pt}{mod-blue-ln, dashed},
    colback=mod-blue-bg
}
\newcounter{daisu}
\NewTColorBox[use counter=daisu]{daisu}{}{
    modern-base,
    borderline={0.4pt}{0pt}{mod-gray-ln},
    colback=mod-gray-bg,
    before upper={\fbox{\textbf{\thedaisu}}\quad} % 必要に応じてタイトルなどにカウンタの値を表示できます
}

% sbox0: 実線・薄い緑背景（公式や重要な定義向け）
\NewTColorBox{sbox0}{}{
    modern-base,
    borderline={0.4pt}{0pt}{mod-green-ln},
    %colback=mod-green-bg
    colback=mod-accent-bg
}

% sbox1: 実線・薄いグレー背景（別のアプローチや参考情報向け）
\NewTColorBox{sbox1}{}{
    modern-base,
    borderline={0.4pt}{0pt}{mod-gray-ln},
    colback=mod-gray-bg
}

% Problem用の共通スタイル
\tcbset{
    modernstyle/.style={
        breakable,
        enhanced,
        frame hidden,
        sharp corners, % あえて角を立てるか、arc=2pt程度で少し丸めるのがモダン
        boxrule=0pt,
        left=10pt,
        right=10pt,
        top=5pt,
        bottom=5pt,
        borderline west={2pt}{0pt}{#1}, % 左側に細い縦線を入れる
    }
}

% Problem環境
\newcounter{prob}
\newtcolorbox[auto counter, use counter=prob]{prob}{
    modernstyle=ProbBorder,
    colback=ProbBlue!50,
    before upper={%
        \textbf{\textcolor{ProbBorder}{Problem \thetcbcounter}} \hspace{0.5em}%
    }
}

% Problem (Section numbering) 環境
\newtcolorbox[auto counter, use counter=prob, number within=section]{prob_s}{
    modernstyle=ProbBorder,
    colback=ProbBlue!50,
    before upper={%
        \textbf{\textcolor{ProbBorder}{Problem \thetcbcounter}} \hspace{0.5em}%
    }
}

% Example環境
\newcounter{exam}
\newtcolorbox[auto counter, use counter=exam, number within=section]{exam_s}{
    modernstyle=ExamBorder,
    colback=ExamGreen!60,
    before upper={%
        \textbf{\textcolor{ExamBorder}{Example \thetcbcounter}} \hspace{0.5em}%
    },
    after upper={\hfill\textcolor{ExamBorder}{\ding{72}}}
}
